\documentclass[11pt,a4paper]{article}
\usepackage[T1]{fontenc}
\usepackage[utf8]{inputenc}
\usepackage{amsmath,amssymb,xcolor}
\usepackage[margin=25mm]{geometry}
\usepackage[colorlinks=true,urlcolor=blue]{hyperref}
\definecolor{nexorablue}{HTML}{2459B5}
\setlength{\parindent}{0pt}
\setlength{\parskip}{7pt}
\setlength{\emergencystretch}{2em}
\renewcommand{\familydefault}{\sfdefault}
\begin{document}

{\small\bfseries\color{nexorablue} NEXORA CPGE}\par
{\small MPSI / NEX-CIR-05 / Corrigé}\par
{\LARGE\bfseries Influence du voltmètre}\par
\vspace{0.25em}\hrule\vspace{0.6em}

\section*{Tension sans appareil}
L’équation $(9-V_B)/60=V_B/30$ donne $V_B=3\,\mathrm V$.

\section*{Tension lue et écart}
La loi des nœuds devient $(9-V_B)/60=V_B/30+V_B/90$. Après multiplication par $180$, $27-3V_B=8V_B$, soit $V_B=27/11\,\mathrm V\simeq2{,}45\,\mathrm V$. L’écart relatif signé vaut $-2/11\simeq-18{,}2\,\%$. Le courant supplémentaire dans le voltmètre explique la baisse.


\end{document}
